\subsection{问题的共同结构}
四个问题问的是同一个物理过程在不同参数体系与不同时间尺度上的表现，因此不应当分别建立四个模型，而应当建立\textbf{一个}带参数的初边值问题，用同一套离散与积分方法求解，四个问题只在下列三处不同：

\begin{center}
\begin{tabular}{@{}llll@{}}
\toprule
问题 & 物性公式 & 时间范围与输出 & 计算域 \\
\midrule
1 & 附录 2（常物性 + $D(C)$） & $0$--\SI{1800}{\second}，每 \SI{1}{\second} & 固定，$r\in[0,R_0]$ \\
2 & 附录 3（$\rho,c_p,k$ 随 $C$；$D(C,T)$） & $0$--\SI{3}{\hour}，每 \SI{1}{\second} & 固定 \\
3 & 附录 3 & $0$--$t_{\mathrm{end}}$，每 \SI{60}{\second} & 固定 \\
4 & 附录 4 & $0$--$t_{\mathrm{end}}$，每 \SI{60}{\second} & 收缩，$r\in[0,R(t)]$ \\
\bottomrule
\end{tabular}
\end{center}

\subsection{关键困难与对策}
\begin{enumerate}[leftmargin=*,itemsep=0.25em]
  \item \textbf{三个量级的时间尺度分离。} 以附录 2 的参数计，热扩散时间尺度 $R_0^2/\alpha=\valQoneHeatTimeScaleS$ s，而水分扩散时间尺度 $R_0^2/D(C_0)=\valQoneMassTimeScaleH$ h，两者相差 $\valQoneStiffnessRatio$ 倍；温度与水分又通过 $D(C,T)$ 强耦合。这使方程组\textbf{刚性}：显式格式的稳定步长由温度方程决定（$\Delta t\lesssim\Delta r^2/2\alpha$，在生产网格上约 \SI{3e-3}{\second}），而需要积分的总时长达 $\valQthreeDrySeconds$ \si{\second}。对策是全隐式、变阶变步长的 BDF 方法（第~\ref{subsec:time} 小节）。
  \item \textbf{扩散系数跨越三到四个数量级。} $D\propto\mathrm{e}^{-a/C}$ 在干燥后期急剧变小：以附录 3 计，$D$ 由初始的 \sci{\valDmaxAppendixThree} \si{\square\meter\per\second} 降到烘干结束时表面的 \sci{\valDsurfaceAtDryQthree} \si{\square\meter\per\second}，比值 $\valDratioQthree$（图~\ref{fig:properties}）。表面因此形成低扩散率的“硬壳”，传质 Biot 数 $h_mR_0/D$ 由 $\valQoneBiotMass$ 增至约 $\valBiotMassAtDryQthree$，过程由\textbf{内部扩散}而非外部对流控制——这一判断在灵敏度分析中被定量证实（弹性 $\valElasH$ 对 $h$、$\valElasDzero$ 对 $D_0$）。
  \item \textbf{边界数据是折线与样条。} 附件 1 每 \SI{60}{\second} 一个样本，线性插值后斜率在每个样本处跳变；附件 2 的 PCHIP 插值在每个节点处曲率跳变。多步法跨过折点时的误差不受容差控制（这一现象在初期实验中被观测到，见 MDR-0004 补充），对策是在\textbf{每个数据折点重启积分器}，段内强迫项光滑。
  \item \textbf{“各处低于 0.15”是一个停止时刻问题。} 判据定义的时刻不落在输出网格上。对策是把 $g(t)=\max_i C_i(t)-0.15$ 作为积分器的事件函数精确求根，并同时报告 \SI{60}{\second} 网格上的值以与结果文件对照。
  \item \textbf{问题 4 是移动边界问题。} 计算域随 $R(t)$ 收缩。若把附录方程视为实验室坐标下的方程（Eulerian），被边界扫过的物质离开计算域，水分总量不守恒；本文以干物质守恒的物质坐标（Lagrangian）为主模型，两者的差异作为解释不确定性的一部分定量给出（第~\ref{subsec:q4model} 小节）。
  \item \textbf{四位小数的精度要求。} 离散误差必须远小于 $5\times10^{-5}$。对策是把输出点取在网格节点上（$\Delta r=0.1/2^k$ \si{\centi\meter}），用收敛表确定 $k$，并用解析级数解独立确认误差量级（第~\ref{sec:validation} 节）。
\end{enumerate}

\subsection{技术路线}
图~\ref{fig:route} 给出总体技术路线。所有计算在同一云端运行（run）内完成，每个阶段的输入输出散列、参数与环境写入清单，论文中的每一个数字都追溯到某个阶段的产物（附录~\ref{app:repro}）。

\begin{figure}[htbp]
\centering
\begin{tikzpicture}[
  font=\small\sffamily,
  node distance=5mm,
  box/.style={draw=ink!70,fill=ink!5,rounded corners=2pt,align=center,inner sep=4pt,minimum height=8mm},
  core/.style={box,fill=accent!10,draw=accent!70},
  chk/.style={box,fill=black!4,draw=black!40},
  ar/.style={-{Latex[length=2mm]},draw=ink!70,thick},
]
\node[box,text width=32mm] (data) {附件 1、附件 2、附录 2--4\\[-1pt]\footnotesize 契约校验（单调性、步长、范围）};
\node[core,right=8mm of data,text width=36mm] (pde) {统一初边值问题\\[-1pt]\footnotesize 守恒律 $\Rightarrow$ 一维径向热质耦合方程\\[-1pt]\footnotesize 对称条件 + Robin 边界};
\node[core,right=8mm of pde,text width=34mm] (num) {顶点型有限体积 $+$\\[-1pt]分段自适应 BDF\\[-1pt]\footnotesize 事件求根定位 $t_{\mathrm{end}}$};
\node[box,below=8mm of data.south west,anchor=north west,text width=20mm] (q1) {问题 1\\\footnotesize 附录 2};
\node[box,right=4mm of q1,text width=20mm] (q2) {问题 2\\\footnotesize 附录 3};
\node[box,right=4mm of q2,text width=20mm] (q3) {问题 3\\\footnotesize 判据 0.15};
\node[box,right=4mm of q3,text width=24mm] (q4) {问题 4\\\footnotesize 附录 4 $+$ Landau};
\node[chk,below=7mm of q1.south west,anchor=north west,text width=52mm] (ver) {\textbf{独立校验（与求解器代码分离）}\\[-1pt]\footnotesize 贝塞尔级数解、网格/时间收敛、守恒量、\\[-2pt]\footnotesize 极值原理、退化检验、二维轴对称交叉检验};
\node[chk,right=5mm of ver,text width=34mm] (sen) {\textbf{灵敏度与备选解释}\\[-1pt]\footnotesize $h,h_m,D_0,C_{\mathrm{air}},T_{\mathrm{air}},s$ 的弹性\\[-2pt]\footnotesize 平台/判据/离散化的备选};
\node[box,below=7mm of ver.south west,anchor=north west,text width=92mm] (out) {结果工作簿 \nolinkurl{result1--4.xlsx}（模板生成、契约检查）\quad$\bullet$\quad 图 / 三线表 \quad$\bullet$\quad 论文与支撑材料（散列清单）};
\draw[ar] (data) -- (pde);
\draw[ar] (pde) -- (num);
\draw[ar] (num.south) -- ++(0,-3mm) -| (q4.north);
\draw[ar] (num.south) -- ++(0,-3mm) -| (q3.north);
\draw[ar] (pde.south) -- ++(0,-3mm) -| (q2.north);
\draw[ar] (data.south) -- ++(0,-3mm) -| (q1.north);
\draw[ar] (q1.south) -- ++(0,-3mm) -| (ver.north);
\draw[ar] (q4.south) -- ++(0,-3mm) -| (sen.north);
\draw[ar] (ver.south) -- ++(0,-3mm) -| (out.north);
\draw[ar] (sen.south) -- ++(0,-3mm) -| ([xshift=-20mm]out.north east);
\end{tikzpicture}
\caption{技术路线：由守恒律建立统一的一维径向热质耦合初边值问题，四个问题只改变物性、时间范围与边界是否移动；每一组解都经与求解器分离的独立校验，并给出灵敏度与备选解释。}
\label{fig:route}
\end{figure}
